to an algebraically closed residue field would become unnecessary, since @@M0@@ is available here as a ``parameter curve'' along which to deform cycles. A difficulty nevertheless arises, for with the notation of 7.12, @@M1@@ must now denote a cycle on @@M2@@ such that @@M3@@ ``is defined,'' but it by no means follows that the cycle @@M4@@ is also defined. As was to be expected, one is thus reduced to formulating a statement of the type of Chow's ``moving lemma'' in order to define (7.14.2).

Observe, however, that the existence of homomorphism (7.4.1) shows that (7.14.1) is compatible with the equivalence relation on cycles defined by the surjective homomorphism @@M5@@, which is only slightly coarser than linear equivalence; thus, when @@M6@@ is quasi-projective and smooth over a field, we point out elsewhere that it differs from linear equivalence only up to torsion (XIV 4.1, 4.2). In this case one therefore obtains a homomorphism analogous to (7.14.2), or better to (7.14.2 bis), but with both members tensored by @@M7@@.

The most natural way to define a homomorphism (7.14.2) directly would seem to be to proceed as in 7.3. This would, however, require a theory of Chow groups @@M8@@ (or of the Chow ring in the regular case) for schemes more general than schemes quasi-projective over a field, including @@M9@@, and a contravariant law for these groups that would allow one to define the restriction homomorphism to @@M10@@, @@M11@@. Compare XIV 8. Once such a sufficiently general theory of Chow groups and rings is available, it should be possible to transcribe into Chow theory all the results enumerated above for the specialization homomorphism in @@M12@@-theory.

